% two-column-ragged — two-column's concern under ragged-right type and spaced paragraphs.
%
% two-column.tex is justified, so a line that wraps runs to the column's edge and that
% is what tells its break from a paragraph's end. Ragged right takes that away: every
% line stops short. What is left is the word the next column starts with, here "column",
% lower case because the sentence carries on. The paragraphs are set flush with space
% between them, as most word processors set them, so no indent marks them either.
\documentclass[twocolumn]{article}
\usepackage{lmodern}
\usepackage[paperwidth=6in,paperheight=2.0in,margin=0.5in]{geometry}
\pagestyle{empty}
\raggedright
\setlength{\parindent}{0pt}
\setlength{\parskip}{6pt}
% No hyphenation: that is the hyphenation fixture's concern.
\hyphenpenalty=10000 \exhyphenpenalty=10000
\begin{document}
The first paragraph is short and sits at the top of the left column.

The second paragraph is long enough to run down the rest of the left column and on
into the right one, so the column break falls in the middle of a sentence, which goes
on after the break as if nothing had happened at all.

The third paragraph starts with a capital and closes the right column.
\end{document}
